% Bagian: first-order-logic
% Bab: axiomatic-deduction
% Subbagian: deduction-theorem

% Verifikasi sifat-sifat keterderivasian yang diperlukan untuk himpunan
% konsisten maksimal dalam bab kelengkapan.

\documentclass[../../../include/open-logic-section]{subfiles}

\begin{document}

\iftag{FOL}
      {\olfileid[id]{fol}{axd}{ded}}
      {\olfileid[id]{pl}{axd}{ded}}

\olsection{Teorema Deduksi}

Seperti telah kita lihat, memberikan !!{derivation}s dalam sistem
aksiomatik itu merepotkan, dan !!{derivation}s semacam itu mungkin sulit
ditemukan. Daripada benar-benar menuliskan daftar panjang !!{formula}s,
umumnya lebih mudah berargumen bahwa !!{derivation}s tersebut ada dengan
memanfaatkan beberapa hasil sederhana. Kita telah membuktikan tiga hasil
semacam itu: \olref[ptn]{prop:reflexivity} menyatakan bahwa kita selalu
dapat menegaskan $\Gamma \Proves !A$ apabila diketahui bahwa $!A \in
\Gamma$. \olref[ptn]{prop:monotonicity} menyatakan bahwa jika $\Gamma
\Proves !A$, maka $\Gamma \cup \{!B\} \Proves !A$ juga. Selanjutnya,
\olref[ptn]{prop:transitivity} menyiratkan bahwa jika $\Gamma \Proves !A$
dan $!A \Proves !B$, maka $\Gamma \Proves !B$. Berikut satu hasil
sederhana lainnya, yaitu versi ``meta'' dari modus ponens:

\begin{prop}
  \ollabel{prop:mp} Jika $\Gamma \Proves !A$ dan $\Gamma \Proves !A \lif
  !B$, maka $\Gamma \Proves !B$.
\end{prop}

\begin{proof}
  Kita mempunyai $\{!A, !A \lif !B\} \Proves !B$:
  \begin{derivation}
    1. & $!A$ & Hip.\\
    2. & $!A \lif !B$ & Hip.\\
    3. & $!B$ & 1, 2, MP
  \end{derivation}
  Menurut \olref[ptn]{prop:transitivity}, $\Gamma \Proves !B$.
\end{proof}

Hasil terpenting yang akan kita gunakan dalam konteks ini adalah teorema
deduksi:

\begin{thm}[Teorema Deduksi]
\ollabel{thm:deduction-thm} $\Gamma \cup \{!A\} \Proves !B$ jika dan
hanya jika $\Gamma \Proves !A \lif !B$.
\end{thm}

\begin{proof}
Arah ``jika'' langsung diperoleh. Jika $\Gamma \Proves !A \lif !B$,
maka menurut \olref[ptn]{prop:monotonicity}, $\Gamma \cup \{!A\}
\Proves !A \lif !B$ juga. Selain itu, menurut
\olref[ptn]{prop:reflexivity}, $\Gamma \cup \{!A\} \Proves !A$. Jadi,
menurut \olref{prop:mp}, $\Gamma \cup \{!A\} \Proves !B$.

Untuk arah ``hanya jika'', kita melakukan induksi pada panjang
!!{derivation} dari $!B$ berdasarkan $\Gamma \cup \{!A\}$.

Sebagai basis induksi, kita membuktikan klaim tersebut untuk setiap
!!{derivation} yang panjangnya~$1$. !!^a{derivation} dari~$!B$
berdasarkan $\Gamma \cup \{!A\}$ yang panjangnya~$1$ hanya terdiri atas
$!B$; dan jika derivasi itu benar, maka $!B$ merupakan anggota
$\Gamma \cup \{!A\}$ atau sebuah aksioma. Jika $!B \in \Gamma$ atau
merupakan aksioma, maka $\Gamma \Proves !B$. Kita juga mempunyai
$\Gamma \Proves !B \lif (!A \lif !B)$ menurut
\olref[prp]{ax:lif1}, dan \olref{prop:mp} memberikan
$\Gamma \Proves !A \lif !B$. Jika $!B \in \{!A\}$, maka
$\Gamma \Proves !A \lif !B$ karena !!{formula} terakhir~$!A \lif !B$
sama dengan $!A \lif !A$, dan kita telah !!{derive} formula itu dalam
\olref[pro]{ex:identity}.

Untuk langkah induksi, andaikan !!a{derivation} dari~$!B$ berdasarkan
$\Gamma \cup \{!A\}$ berakhir dengan suatu langkah~$!B$ yang
dijustifikasi oleh modus ponens. Jika langkah itu tidak dijustifikasi
oleh modus ponens\iftag{FOL}{ ataupun aturan~\QR}{}, maka
$!B \in \Gamma$, $!B \ident !A$, atau $!B$ merupakan aksioma, dan
penalaran yang sama seperti pada basis induksi berlaku.
\iftag{FOL}{Kasus aturan~\QR{} ditangani dalam bagian berikutnya.}{}
Dengan demikian, sejumlah
langkah sebelumnya dalam
!!{derivation} tersebut adalah $!C \lif !B$ dan $!C$, untuk suatu
!!{formula}~$!C$;
yakni, $\Gamma \cup \{!A\} \Proves !C \lif !B$ dan
$\Gamma \cup \{!A\} \Proves !C$. !!^{derivation}s masing-masing lebih
pendek, sehingga hipotesis induksi berlaku untuk keduanya. Jadi, kita
mempunyai kedua hal berikut:
\begin{align*}
  & \Gamma \Proves !A \lif (!C \lif !B); \\
  & \Gamma \Proves !A \lif !C.
\end{align*}
Namun, menurut \olref[prp]{ax:lif2}, kita juga mempunyai
\[
\Gamma \Proves (!A \lif (!C \lif !B)) \lif
((!A\lif !C)  \lif (!A \lif !B)),
\]
dan dua penerapan \olref{prop:mp} memberikan
$\Gamma \Proves !A \lif !B$, sebagaimana diperlukan.
\end{proof}

Perhatikan bagaimana \olref[prp]{ax:lif1} dan \olref[prp]{ax:lif2}
dipilih tepat agar Teorema Deduksi berlaku.

Berikut beberapa fakta berguna tentang !!{derivability} yang kita
tinggalkan sebagai latihan.

\begin{prop}
  \ollabel{prop:derivfacts}
\begin{enumerate}
\item $\Proves (!A \lif !B) \lif ((!B \lif !C)
  \lif (!A \lif !C))$; \ollabel{derivfacts:a}
\item Jika $\Gamma \cup \{ \lnot !A\}
  \Proves \lnot !B$, maka $\Gamma \cup \{ !B\} \Proves
  !A$ (Kontraposisi); \ollabel{derivfacts:b}
\item $\{ !A, \lnot!A\} \Proves
    !B$ (Ex Falso Quodlibet, Eksplosi); \ollabel{derivfacts:c}
\item $\{ \lnot\lnot!A\} \Proves
  !A$ (Eliminasi Negasi Ganda);\ollabel{derivfacts:d}
\item Jika $\Gamma \Proves \lnot\lnot!A$, maka $\Gamma \Proves
  !A$;\ollabel{derivfacts:e}
\end{enumerate}
\end{prop}

\tagprob{FOL}
\begin{prob}
  Buktikan \olref[fol][axd][ded]{prop:derivfacts}.
\end{prob}
\tagendprob

\tagprob{notFOL}
\begin{prob}
  Buktikan \olref[pl][axd][ded]{prop:derivfacts}.
\end{prob}
\tagendprob

\end{document}
